% Planned223.5: explicit proof-version survey and separate original studies.
\chapter{Insight, Impasse and Restructuring}
\label{sec:p19-insight-impasse-restructuring}
\chaptermark{Insight and Restructuring}

Insight research separates how a solution becomes available, how it is
experienced, and whether it satisfies the problem. Sudden confidence is
evidence about experience, not proof of correctness or identification
of one neural module.

The survey source used here is the complete author-hosted
\emph{uncorrected revises proof} of van Steenburgh, Fleck, Beeman and
Kounios's \emph{Insight}, printed pp.~475--491
\parencite{VanSteenburghEtAl2012InsightRevises}.
Its explicit OUP production banner identifies a production proof rather
than the final published 2012 Oxford chapter~24. Direct original evidence comes from
Knoblich and colleagues' 1999 matchstick experiments, Jung-Beeman and
colleagues' 2004 fMRI/EEG study, and Subramaniam and colleagues' 2009
affect/fMRI study \parencite{KnoblichEtAl1999ConstraintsChunks,
JungBeemanEtAl2004NeuralInsight,SubramaniamEtAl2009PositiveAffect}.
This finite corpus is neither an exhaustive replication review nor a
meta-analysis. Dated physiological interpretations are not clinical
or creativity-enhancement prescriptions.

\section{Experience, process and correctness}
\label{sec:p19-insight-operations}

\begin{hbdefinition}{def:p19-insight-distinct-claims}{Four claims that should not be collapsed}
Phenomenological insight is reported sudden awareness of an interpretation
or candidate, often with surprise and confidence.
Operational insight is the study's actual recording category, such as
a postsolution insight/noninsight judgment.
A representation specifies the interpreted elements, goals, allowable
operations and constraints; representational change revises some of that
specification.
Correctness is satisfaction of an independently stated task benchmark.

A retrospective label does not trace every intermediate operation.
A correct answer may be remembered, guessed, derived deliberately or
experienced as sudden. A changed representation can remain inadequate.
Confidence is not the correctness test.
\end{hbdefinition}

The survey treats suddenness, apparent obviousness and the emotional Aha
as typical while acknowledging dispute about their necessity.
Some definitions add a restructuring requirement. Its analytic contrast
involves deliberately manipulating problem elements, not a proof that
all analytic processing is consciously accessible
\parencite[pp.~475--479]{VanSteenburghEtAl2012InsightRevises}.

The category belongs to a solution episode, not permanently to a stimulus.
Classic insight problems can be solved analytically, and ostensibly
analytic tasks can yield sudden realizations. Comparing heterogeneous
puzzle classes confounds strategy with modality, complexity, memory
load and prior knowledge. Trial-level reports within a relatively
homogeneous family reduce these stimulus differences but depend on
interpretation and reporting thresholds
\parencite[pp.~486--488]{VanSteenburghEtAl2012InsightRevises}.

\begin{hbexample}{def:p19-insight-cra}{Compound remote associates and the actual benchmark}
CRAB, PINE and SAUCE admit APPLE: crabapple, pineapple, apple sauce.
TOOTH, POTATO and HEART admit SWEET: sweet tooth, sweet potato, sweetheart.
One word must fit all three; three separately plausible associates do not
solve the task. ACHE fits tooth and heart but not potato.
The compound order can differ across pairings.
Language-specific lexical knowledge and candidate selection therefore matter
\parencite[pp.~501, 507--508]{JungBeemanEtAl2004NeuralInsight}
\parencite[pp.~418--419]{SubramaniamEtAl2009PositiveAffect}.
\end{hbexample}

Generating an answer differs from suddenly recognizing an experimenter-supplied
answer. The survey's x-ray and seesaw cues change comprehension of initially
puzzling sentences; their neural contrasts are not interchangeable with
unaided CRA generation \parencite[pp.~482, 488]{VanSteenburghEtAl2012InsightRevises}.
Nor does a participant's Aha make an incorrect candidate correct.

\section{Prior knowledge, fixation and worked problems}
\label{sec:p19-insight-knowledge}

Knowledge is necessary for recognizing numerals, interpreting equations,
retrieving word meanings and understanding objects. It can also impose
restrictions not demanded by the current task. The historical
reproductive/productive distinction concerns reuse versus reorganization,
not a requirement that useful insight be independent of experience
\parencite[pp.~479--481]{VanSteenburghEtAl2012InsightRevises}.

In the candle problem, a box holding tacks can also become a support after
being emptied and attached. Functional fixedness is the proposed persistence
of the container interpretation. Incremental-learning accounts offer another
route: a failed attempt to push tacks through the candle may redirect attention
to thinner material. The final support idea alone does not distinguish release
from fixedness from learning through unsuccessful attempts.

The radiation/convergence task likewise distinguishes possessing relevant
knowledge from applying it. Several low-intensity paths must converge at the
target without excessive intensity along healthy tissue.
The original transfer evidence is developed separately in
\cref{sec:p19-analogy-and-problem-transfer}.
Its puzzle stipulations are not clinical treatment advice.

The nine-dot task asks for four straight connected segments through a
three-by-three grid without lifting the pencil or retracing a segment.
It does not require turns at dots or confinement within their bounding square.
Nevertheless knowing that outside lines are permitted differs from constructing
the solution. The survey contrasts knowledge, visuospatial construction and
training/hint effects; it does not assign every failure to one implicit
boundary assumption \parencite[pp.~479--480]{VanSteenburghEtAl2012InsightRevises}.

\begin{hbexample}{def:p19-insight-nine-dot}{A candidate whose correctness is separately checkable}
For dots at \(\{0,1,2\}\times\{0,1,2\}\), draw successively through
\[
 (0,0),\ (2,2),\ (2,-1),\ (-1,2),\ (2,2).
\]
The outside turns are legal under the stated instruction.
This coordinate reconstruction is handbook-owned, not a participant record
or a source-exact copy of the survey figure.
\end{hbexample}

\begin{hblemma}{thm:p19-insight-nine-dot-path}{proved=annex:p19.insight-nine-dot}{Correctness of the four-segment construction}
The displayed segments join end to end, contain all nine grid dots and
do not retrace a segment. This verifies the candidate under the stated
benchmark, not its cognitive origin or a minimum-line claim for another task.
\end{hblemma}

\section{Representational change needs a mechanism}
\label{sec:p19-insight-representation-change}

Knoblich and colleagues specify constraint relaxation and chunk decomposition.
The first changes what outcomes or moves are treated as admissible;
the second changes the granularity of perceived components.
Their theories predict relative difficulty and transfer, rather than naming
every eventual success restructuring
\parencite[pp.~1534--1538]{KnoblichEtAl1999ConstraintsChunks}.

A false Roman-numeral arithmetic statement must become true by moving exactly
one matchstick. Moving, sliding and rotation are allowed; discarding is not.
Slanted sticks may not become division signs, and V/X have the prescribed
two-slanted-stick construction. Ordinary algebra's compensating-change rules
are not this task's benchmark.

\begin{tabularx}{\textwidth}{@{}l X X@{}}
\toprule Type & Worked move & Proposed demand\\ \midrule
A & \(\mathrm{VI}=\mathrm{VII}+\mathrm I\) becomes
\(\mathrm{VII}=\mathrm{VI}+\mathrm I\), moving one I. &
Relax value preservation; decompose loose composite numerals.\\
B & \(\mathrm I=\mathrm{II}+\mathrm{II}\) becomes
\(\mathrm I=\mathrm{III}-\mathrm{II}\), moving the vertical plus stick
into the first II. & Relax value/operator constraints; decompose a sign.\\
C & \(\mathrm{III}=\mathrm{III}+\mathrm{III}\) becomes
\(\mathrm{III}=\mathrm{III}=\mathrm{III}\), rotating the vertical plus stick. &
Relax operator and whole-equation-form constraints.\\
D & \(\mathrm{XI}=\mathrm{III}+\mathrm{III}\) becomes
\(\mathrm{VI}=\mathrm{III}+\mathrm{III}\), sliding a component of X to form V. &
Relax value preservation; decompose tight numeral chunks.\\
\bottomrule
\end{tabularx}

VII is loose in the source's sense because V and I remain meaningful units.
V/X are tight because their separate diagonal strokes have no numeral meaning.
Plus/equals are intermediate: their components can be meaningful but are
not commonly parsed independently. The source's
\enquote{tautology constraint} concerns avoiding identity-chain outcomes,
not a logical theorem that every \(X=Y=Z\) string is valid.

The predicted first-encounter time ordering is B slower than A, C slower than B,
and D slower than A. No interval scale supplied a general prediction for D
relative to B/C. The paper explicitly acknowledges that it did not independently
observe constraint-encoded knowledge; the mechanisms remain theoretical
explanations tested through predictions
\parencite[pp.~1535--1538, 1553]{KnoblichEtAl1999ConstraintsChunks}.

\section{Original matchstick evidence and its exceptions}
\label{sec:p19-insight-matchstick-evidence}

The four experiments are named 1A, 1B, 2 and 3.
Participants and repeated problem responses are different units.

\begin{longtable}{@{}lp{0.31\textwidth}p{0.46\textwidth}@{}}
\caption{Original matchstick evidence and its exceptions}\\
\toprule Experiment & Sample/recruitment & Design\\ \midrule
\endfirsthead
\caption[]{Original matchstick evidence and its exceptions (continued)}\\
\toprule Experiment & Sample/recruitment & Design\\ \midrule
\endhead
\midrule\multicolumn{3}{r}{\emph{continued on next page}}\\
\endfoot
\endlastfoot
1A &20 Hamburg undergraduates, 13 women/7 men, ages 19--45,
course credit, Roman-numeral familiarity required.&
Within-person types, two randomized six-problem blocks:
two A, two B, one C, one D each.\\
1B &38 paid participants, 17 men/21 women, ages 20--36,
Munich campus/newspaper recruitment.&
Random assignment to two new stimulus versions; one block.\\
2 &78 Chicago undergraduates, 45 women/33 men, ages 18--38, course credit.&
Random assignment to four numeral/order-similarity conditions;
type/block repeated within people.\\
3 &34 Hamburg/Munich undergraduates, 18 women/16 men, ages 19--42,
course credit.&
Random order groups: 18 regular-first and 16 identity-chain-first.\\
\bottomrule
\end{longtable}

Numeral recognition was trained first: in 1A, five 75-trial blocks with I--XV.
A problem remained until a correct announced move or five minutes.
Unsolved trials were followed by disclosure of the answer.
Consequently later-block improvement includes exposure to taught solutions,
not only spontaneous insight. Randomization within blocks does not remove
all differences among symbolic types
\parencite[pp.~1539,1541,1544,1548]{KnoblichEtAl1999ConstraintsChunks}.

The outcomes were five-minute solution frequency and solution time.
Unsolved trials were entered at 300 seconds; these are capped observations,
not uncensored eventual latencies. Skewed distributions motivated reported
medians/quartiles and ordinal tests.

\begin{tabular}{@{}lrrrr@{}}
\toprule 1A capped time (seconds) &A&B&C&D\\ \midrule
Block 1 median &40&155&300&140\\
Block 1 lower quartile &14&67&91&98\\
Block 1 upper quartile &61&190&300&290\\
Block 2 median &33&38&10&42\\
\bottomrule
\end{tabular}

The 1A solution-frequency type-by-block interaction was
\(F(3,57)=3.24,p<.05\). D's frequency increased 75\% to 90\% without
a significant transfer test (\(p=.13\)), while its capped-time transfer
test was significant (\(p<.05\)). These are different outcomes
\parencite[pp.~1540--1541,tables~2--3]{KnoblichEtAl1999ConstraintsChunks}.

In 1B, medians were 43, 126, 300, 104 seconds for A, B, C, D.
A--B frequency separation was inconclusive (\(p=.12\)), although its time
contrast was significant. Ceiling effects are the authors' explanation,
not independently confirmed from rounded summaries.
New stimulus versions replicate parts of a within-paper pattern,
not every problem requiring insight.

Experiment 2 groups had 19 same/same, 20 same-numerals, 19 same-order,
20 different/different participants. Rounded solution percentages rose
respectively 67 to 96, 64 to 99, 64 to 95, 71 to 95 across blocks.
The reported similarity frequency test was \(F(3,74)=.34,p=.79\):
not proof of a zero similarity effect. C's time analysis did find
\(H(3,N=78)=11.05,p<.05\), while D's test gave \(p=.08\).
Later C/D times also became faster than A/B rather than merely equal,
a reversal acknowledged with post hoc explanations
\parencite[pp.~1544--1547,tables~7--9]{KnoblichEtAl1999ConstraintsChunks}.

Experiment 3 reversed C's move: an initially false identity chain became
a conventional calculation by turning equals into plus.
First-encounter solution rates were 61\% for regular C and 100\% for
identity-chain C, with reported \(\chi^2(1,N=34)=12.59,p<.001\).
For D, 61\% versus 69\% was inconclusive (\(p=.64\)).
The regular C median was 175 seconds on first encounter versus 11 seconds
after prior identity-chain problems; prior exposure and block also differ.
This supports discrimination of demands, not direct observation of a
private constraint being removed
\parencite[pp.~1548--1550,tables~10--12]{KnoblichEtAl1999ConstraintsChunks}.

Generic reorganization, variation-and-test, functional fixedness, mental ruts,
association hierarchies and procedural recall need additional assumptions
to predict these particular difficulty/transfer patterns.
The authors' criticism of underspecified alternatives does not prove that
every possible implementation is impossible. Unchanged displays and no
imposed break challenge a mechanism requiring all impasses to end only
through external change or time away, not every incubation account
\parencite[pp.~1550--1554]{KnoblichEtAl1999ConstraintsChunks}.

The record retains unfavorable and unexpected results and supplies no
new confidence interval, participant-level reanalysis, causal mediation
estimate or population prevalence. Printed statistical prose sometimes
reverses type comparisons; numerical tables and stated predictions govern
the directional description. Sample frames, answer disclosure, item sets
and the ceiling limit generalization.

\section{Impasse, incubation and competing accounts}
\label{sec:p19-insight-impasse-incubation}

Impasse can mean reported inability to generate another approach, a coded
interval without progress, repeated failed moves, or a gaze proxy.
Each needs thresholds and a recording rule. Eye fixation is not cognitive
fixation by definition: longer gaze may reflect inspection, difficulty,
fatigue or strategy. Verbal protocols can omit inaccessible processing or
alter it through the instruction to speak.

Representational-change theory explains difficulty through an inadequate
initial space. Progress monitoring permits strategy revision when expected
progress is lacking, without requiring a prior felt impasse.
Incremental learning permits a sudden answer report after gradual acquisition.
Sudden awareness and gradual antecedents are therefore compatible
\parencite[pp.~481--483,487]{VanSteenburghEtAl2012InsightRevises}.

The survey's number-reduction illustration uses explicit discovery of a
regularity: its ABCCB response structure makes the last response repeat the
second, permitting a shortcut. Knowledge of that shortcut, a warmth judgment
and an ERP are separate measurements; no sample or effect size is inferred
from the illustration.

Incubation is an interval away from direct work, not a synonym for impasse
or insight. Selective forgetting, continued unconscious work, attention
switching, recovery from fixation and encountering external hints can all
change during a break. A benefit alone does not choose among them.
The survey contrasts fully removed versus partly visible problems, misleading
initial material versus its absence, and demanding intervening tasks versus
undirected rest \parencite[pp.~483--484]{VanSteenburghEtAl2012InsightRevises}.

A discriminating design would measure the proposed impasse before assigning
breaks, equate total exposure/test opportunities, independently vary
misleading primes and relevant external information, and include continuous-work
and time-matched controls. This is a methodological proposal, not a completed
study or a general instruction to take breaks.
This direct original-study corpus supplies no new incubation effect estimate
or meta-analytic average; original incubation papers retain separate
evidence requirements when their numerical results are used.

\section{Original neural evidence: one problem family, separate samples}
\label{sec:p19-insight-neural-evidence}

Jung-Beeman and colleagues categorized correct CRA solutions by participants'
retrospective reports rather than assigning heterogeneous puzzle classes.
Eighteen people entered fMRI testing (10 men, 8 women); four men and one
woman were excluded for poor signal or fewer than ten insight/noninsight
responses, leaving 13. A separate EEG sample had 19 new people (9 men, 10 women).
They were paid, neurologically intact, right-handed native English speakers
aged 18--29 \parencite[pp.~507--508]{JungBeemanEtAl2004NeuralInsight}.

The fMRI task used 124 items and a 30-second limit. A bimanual press indicated
solution, followed by a spoken answer and an insight/no-insight judgment.
An unrelated line task separated trials. EEG used similar procedures with
smaller displays, inter-event delays and no intervening line task.
The principal contrast selected correct insight versus correct noninsight
solutions.

The fMRI sample solved 59\% of items and classified approximately 56\% of
solutions insight and 41\% noninsight.
Mean solution times were 10.25 versus 11.28seconds, SD 3.58 versus 4.13;
the difference was inconclusive (\(t<1,p>.3\)).
A right anterior superior temporal gyrus cluster exceeded the 500-\(\mathrm{mm}^3\)
criterion at voxel threshold \(t(12)>3.43,p<.005\) uncorrected.
The peak voxel difference was 0.25\% across the analysis window;
12 of 13 showed the positive contrast and one the opposite.
A selected peak is not every region's or participant's average effect
\parencite[pp.~501--503,508]{JungBeemanEtAl2004NeuralInsight}.

The EEG sample solved 46\% (SD 8.2), labeling 56\% of correct solutions
insight (SD 8.4) and 42\% noninsight (SD 9.0).
The 39-Hz gamma contrast appeared about 0.3seconds before the solution press.
The insight/time/hemisphere interaction was \(F(1,18)=8.11,p=.011\);
the right T8 contrast in the \(-.30\) to \(-.02\)-second interval was
\(t(18)=3.48,p=.003\).
Planned theta/beta tests were not significant.
Posterior alpha differences were reported; interpreting them as sensory
gating and gamma as conscious semantic integration requires additional
physiological/timing assumptions
\parencite[pp.~504--508]{JungBeemanEtAl2004NeuralInsight}.

Shared response sequences and the bimanual press address some motor alternatives.
They do not randomize people into cognitive mechanisms or trace every step.
Scalp source localization is model-based; BOLD is delayed hemodynamics.
Other regions' subthreshold effects are not proof of no contribution.
Both styles recruited shared networks. The proposed right-temporal integration
account was not a lesion/perturbation necessity or sufficiency test.
Read it with \cref{def:p19-ncm-claims}, not as a diagnostic neural signature.
Selecting correct solutions also prevents this contrast from estimating
the error rate of all Aha responses.

\section{Original affect evidence: assessed state is not an intervention}
\label{sec:p19-insight-affect-evidence}

Subramaniam and colleagues assessed PANAS affect, STAI anxiety and personality
before 135 CRA items; they did not randomly induce positive mood.
They report 52 behavioral-only participants and 30 scanned, with three
imaging exclusions (two poor signals; one with only two analytic responses).
The retained dataset was 79, including 27 for fMRI.
All were neurologically healthy, right-handed native English speakers.
This retained pool is not a complete recruitment account
\parencite[pp.~418--420]{SubramaniamEtAl2009PositiveAffect}.

The 15-second limit differs from the 2004 study's 30 seconds.
Fixation/inter-event delays were jittered zero to eight seconds.
PA minus NA was the positive-mood index; PA/NA need not be opposite poles.
Only five of 79 had NA above 20 on its 10--50 scale, limiting its represented range.
High/low-thirds behavioral subgroups had 26 people each.
Solving percentages use 135 items as denominator; insight percentages use
solved items.

\begin{tabular}{@{}lrrrr@{}}
\toprule Group &Mean solved &Mean insight count &Solved \% &Insight \%\\
\midrule
All 79 &55.3&28.1&41.0&50.8\\
High PA--NA &60.0&34.5&44.4&57.5\\
Low PA--NA &51.3&21.9&38.0&42.7\\
High anxiety &52.3&24.1&38.7&46.1\\
Low anxiety &57.1&33.1&42.3&57.9\\
\bottomrule
\end{tabular}

The high/low mood solved-count comparison was \(t(50)=2.24,p<.05\);
the insight-count comparison was \(t(50)=3.96,p<.0005\).
Analytic counts did not reliably differ across the three mood levels
(\(p=.233\)).
The reported regression including mood/personality measures related PA--NA
to insight percentage with \(r=.40\), Results \(p<.005\), and anxiety
with \(r=-.34,p<.005\).
Anxiety's overall-solving group test was inconclusive (\(p=.207\));
a larger analytic proportion (\(p=.033\)) did not establish a larger
raw analytic count (\(p=.222\))
\parencite[pp.~421--422,426]{SubramaniamEtAl2009PositiveAffect}.

More overall solutions makes pure relabeling inadequate by itself,
but these observational comparisons do not show that inducing mood causes
the observed difference. Ability, engagement, prior knowledge, anxiety and
unmeasured factors remain possible contributors.
Rounded means/percentages do not license raw counts or new confidence intervals.

Across 27 imaging participants, preparatory dACC responsivity correlated
with PA--NA, reported \(r(25)=.41,p<.05\).
ROI definitions included preparation-responsive regions, insight-sensitive
regions and high/low mood differences. The extreme imaging groups had
eight people each because tied scores prevented nine-per-tail thirds.
Convergence emphasized dACC, but not every candidate region shared the mood
relation. The authors consider detecting competing responses, switching
attention, broader access, reduced anxiety and default-state modulation,
and explicitly state that the precise ACC control operation was not isolated
\parencite[pp.~420--430,note~1]{SubramaniamEtAl2009PositiveAffect}.

Preparatory timing constrains an account; it is not proof of a mood--region--insight
causal mediation chain or an insight module.
Overlapping functional contrasts are not independent replications.
The reported tests/SDs retain their own uncertainty meaning; no new
participant-level fit or confidence interval is supplied.

\section{Transfer, training and the remaining theoretical questions}
\label{sec:p19-insight-alternatives}

Representational change describes what becomes admissible or decomposable.
Progress monitoring describes when strategy is revised.
Incremental learning describes information acquired through failure;
remote-association accounts describe availability of candidates.
Forgetting and external-cue accounts describe distinct exits from an
unproductive candidate. They can interact. One final successful move does
not uniquely identify the mechanism.

The survey favors special processes over a strong claim of no difference
whatsoever, yet its future directions allow gradual unconscious accumulation
before sudden access and demand study of sequencing/connectivity, not only
localization. Its warmth, verbalization, working-memory and gaze examples
remain attributed survey reports rather than new complete original-study
records in this corpus \parencite[pp.~478--483,486--488]{VanSteenburghEtAl2012InsightRevises}.

Working memory may support search in a large initial space without explaining
revision isolated in a small one. Puzzle-specific training is not a demonstrated
domain-general insight capacity. Gaze, attention and movement manipulations
can provide relevant information rather than reveal spontaneous insight.
Preparatory and resting-state contrasts differ from perturbations.
The survey's brain-stimulation forecasts remain dated proposals, not
established enhancement procedures.

The German/US matchstick comparison and additional stimulus versions give
limited within-paper generalization evidence.
The 2009 CRA study changed sample, scanner, limit and question; it is not an
exact replication of all 2004 EEG effects. The corpus supplies no effect prevalence
across languages, age groups, clinical populations or scientific discoveries.
Generating versus recognizing an answer, verbal versus perceptual tasks,
and short puzzles versus sustained creativity remain separate transfer tests.

\begin{hbmethod}{def:p19-insight-evidence-comparison}{An inspectable comparison}
State instructions, legal moves and correctness first.
Record prior exposure/domain knowledge.
Keep answer success, insight label, reported impasse, coded representation
change and neural measure separate.

Specify a discriminating mechanism prediction: initial difficulty by constraint
type, transfer after changed form, progress before reported insight, or effects
of misleading/neutral material under independently assigned breaks.
Report assignment, participant counts, repeated trials, exclusions, time limits,
answer disclosure, unfavorable and inconclusive outcomes.

For neural claims identify task phase, contrast and spatial/time/statistical
model. Correlation is not necessity. For facilitation separate assessed from
induced affect and solving rate from insight proportion.
Use \cref{def:p19-evidence-record,def:p19-uncertainty-terms}.
Then check the candidate against the actual benchmark:
the geometric path, arithmetic statement, or all three CRA pairings.
Neither certainty nor a brain signal replaces that check.
This is an editorial evidence method, not a validated universal recipe for insight.
\end{hbmethod}

\begin{proofboundary}{def:p19-insight-source-boundaries}{Source-version and empirical limits}
The survey's revises-proof version and the three original study records
retain separate identities. The reported designs and aggregate results do
not supply a participant-level reanalysis, a causal mood--region--insight
mediation estimate, or a general prescription for improving insight.
The production proof is not identified with the final published chapter.
\end{proofboundary}

\index{insight}\index{impasse}\index{restructuring}
\index{constraint relaxation}\index{chunk decomposition}
\index{incubation}\index{Aha experience}
